% ---------------------------------------------------------------------------
% Supplemental Material: details behind Sec. VII (Results) of the main text
% ---------------------------------------------------------------------------

This section gives the full benchmark tables behind Sec.~\ref{sec:results} and the
details the main text summarises. Numbers quoted for a ring family without its overlap
level are at $o_n=0.05$.

\emph{Recovery window and fusions.} At $o_n=0.05$, for $k=16$ the isolation hypothesis
of Theorem~\ref{thm:window} cannot hold, and for $k=8$ it would need exactly five merge
positions more than $k$ apart on every ring. Rings are then fused inside the window in
three ways: two rings share two merged nodes within ring distance $k$ on both rings
(Proposition~\ref{prop:sharedpairs}); they share two merged nodes within $k/2$ on one
ring and within $3k/2$ on the other; or the merged nodes of three pairwise-merged rings
form an admissible wedge. Without noise, between $6\%$ and $61\%$ of the rings of a family
are fused at $o_n=0.05$ and between $38\%$ and all of them at $o_n=0.10$
(Table~\ref{tab:doublebonds}). For $k=16$ the noise-free curve of
Fig.~\ref{fig:omegat}(b) reaches $0.95$ only for $t\ge0.30$.

\begin{table*}[t]
\caption{Accuracy on the benchmarks as a user would run the methods, every parameter chosen without labels by one rule: on each test graph, the parameter whose cover has the largest overlapping modularity, Eq.~\eqref{eq:qov}, which for a partition is the Newman--Girvan modularity (label-free protocol). $\Omega$, mean over the noise range ($p\in[0,0.8]$ for rings, $\mu\in[0.1,0.6]$ for LFR), ten test graphs per noise level. HC: HyperCommon; e-c.: its triples joined only through shared edges; tri.: triangles only (J-thresholded 3-clique percolation~\cite{palla2005uncovering,kumpula2008sequential}); 2-test: third pair not tested; Link: link communities with the same similarity; CPM: clique percolation; Leiden: Leiden with the Reichardt--Bornholdt modularity; L-CPM: Leiden with the Constant Potts Model; Info.: Infomap; Walk.: Walktrap, walks of length 4, dendrogram cut at maximum modularity; LPA: label propagation; SBM: nested degree-corrected stochastic block model, each level of its hierarchy a parameter value; OSLOM: order statistics local optimisation, significance threshold $t$ as its parameter; SCAN: SCAN with cosine similarity; SCAN-J: SCAN with Jaccard similarity; $\alpha$: exponent of the preferential (degree-biased) noise. Parameter-free methods (Walktrap, LPA) are as under every protocol. Superscripts on HC: the paired 95\% bootstrap interval of its difference from the best other method (ablations excluded) lies above zero ($+$), contains zero ($=$) or lies below zero ($-$). Bold: best in row, within 0.005. 95\% bootstrap half-widths of these means over the ten test seeds have median 0.007 and maximum 0.050.}
\label{tab:labelfree}
\scriptsize\setlength{\tabcolsep}{1.35pt}
\resizebox{\textwidth}{!}{%
\begin{tabular}{llcccc|cc|cccccc|cccccc}
\hline\hline
 & & \multicolumn{4}{c}{HC and ablations} & \multicolumn{2}{c}{relatives} &\multicolumn{6}{c}{partitions} & \multicolumn{6}{c}{overlapping methods} \\
Family & $o_n$ & HC & e-c. & tri. & 2-test & Link & CPM & Leiden & L-CPM & Info. & Walk. & LPA & SBM & SLPA & OSLOM & DEMON & ANGEL & SCAN & SCAN-J \\
\hline
uniform, $k{=}8$ & 0.00 & \textbf{0.61$^{=}$} & \textbf{0.61} & 0.51 & 0.51 & 0.60 & 0.51 & \textbf{0.61} & \textbf{0.61} & 0.48 & 0.39 & 0.14 & 0.38 & 0.36 & 0.55 & 0.53 & 0.52 & 0.50 & 0.53 \\
 & 0.05 & \textbf{0.55$^{+}$} & 0.54 & 0.49 & 0.29 & 0.53 & 0.49 & 0.44 & 0.43 & 0.32 & 0.24 & 0.12 & 0.27 & 0.33 & 0.45 & 0.44 & 0.39 & 0.26 & 0.30 \\
 & 0.10 & \textbf{0.50$^{+}$} & \textbf{0.50} & 0.46 & 0.22 & 0.45 & 0.44 & 0.32 & 0.31 & 0.25 & 0.18 & 0.11 & 0.21 & 0.30 & 0.37 & 0.40 & 0.40 & 0.21 & 0.21 \\
\hline
uniform, $k{=}16$ & 0.00 & 0.78$^{-}$ & 0.78 & 0.75 & 0.69 & 0.76 & 0.67 & \textbf{0.84} & 0.84 & 0.47 & 0.65 & 0.14 & 0.50 & 0.34 & 0.72 & 0.62 & 0.67 & 0.65 & 0.59 \\
 & 0.05 & 0.72$^{-}$ & 0.71 & 0.69 & 0.39 & 0.54 & 0.56 & \textbf{0.75} & 0.75 & 0.35 & 0.47 & 0.11 & 0.38 & 0.32 & 0.64 & 0.50 & 0.63 & 0.33 & 0.35 \\
 & 0.10 & \textbf{0.62$^{=}$} & \textbf{0.62} & 0.60 & 0.29 & 0.39 & 0.54 & \textbf{0.62} & 0.61 & 0.29 & 0.33 & 0.09 & 0.28 & 0.29 & 0.54 & 0.45 & 0.57 & 0.24 & 0.25 \\
\hline
equal $k$, unequal $s$, $k{=}8$ & 0.00 & \textbf{0.57$^{+}$} & 0.57 & 0.46 & 0.49 & 0.56 & 0.46 & 0.56 & 0.56 & 0.36 & 0.33 & 0.10 & 0.33 & 0.25 & 0.47 & 0.48 & 0.49 & 0.47 & 0.49 \\
 & 0.05 & \textbf{0.52$^{+}$} & 0.51 & 0.45 & 0.25 & 0.49 & 0.43 & 0.35 & 0.34 & 0.24 & 0.19 & 0.09 & 0.23 & 0.23 & 0.34 & 0.38 & 0.35 & 0.22 & 0.23 \\
 & 0.10 & \textbf{0.45$^{+}$} & \textbf{0.45} & 0.40 & 0.17 & 0.41 & 0.37 & 0.24 & 0.24 & 0.19 & 0.15 & 0.08 & 0.16 & 0.21 & 0.27 & 0.33 & 0.33 & 0.15 & 0.16 \\
\hline
equal $k$, unequal $s$, $k{=}16$ & 0.00 & 0.80$^{-}$ & 0.79 & 0.77 & 0.70 & 0.77 & 0.71 & \textbf{0.84} & \textbf{0.84} & 0.60 & 0.69 & 0.17 & 0.49 & 0.43 & 0.76 & 0.64 & 0.69 & 0.66 & 0.60 \\
 & 0.05 & \textbf{0.73$^{=}$} & \textbf{0.73} & 0.71 & 0.42 & 0.56 & 0.59 & \textbf{0.73} & 0.72 & 0.51 & 0.56 & 0.14 & 0.40 & 0.41 & 0.70 & 0.51 & 0.66 & 0.37 & 0.37 \\
 & 0.10 & \textbf{0.65$^{=}$} & \textbf{0.65} & 0.62 & 0.31 & 0.42 & 0.57 & 0.60 & 0.59 & 0.44 & 0.44 & 0.12 & 0.34 & 0.37 & 0.64 & 0.45 & 0.61 & 0.27 & 0.29 \\
\hline
equal $s$, unequal $k$ & 0.00 & 0.78$^{=}$ & \textbf{0.78} & 0.76 & 0.71 & 0.74 & 0.60 & \textbf{0.79} & \textbf{0.79} & 0.55 & 0.60 & 0.16 & 0.51 & 0.35 & 0.56 & 0.53 & 0.64 & 0.66 & 0.61 \\
 & 0.05 & 0.70$^{=}$ & \textbf{0.70} & 0.69 & 0.47 & 0.57 & 0.59 & 0.67 & \textbf{0.70} & 0.44 & 0.52 & 0.13 & 0.40 & 0.31 & 0.50 & 0.38 & 0.57 & 0.40 & 0.41 \\
 & 0.10 & 0.61$^{-}$ & 0.61 & 0.59 & 0.39 & 0.45 & 0.56 & 0.58 & \textbf{0.63} & 0.39 & 0.46 & 0.11 & 0.34 & 0.27 & 0.46 & 0.29 & 0.46 & 0.32 & 0.34 \\
\hline
two worlds & 0.00 & \textbf{0.83$^{+}$} & \textbf{0.82} & 0.82 & 0.71 & 0.78 & 0.70 & 0.58 & 0.79 & 0.40 & 0.52 & 0.13 & 0.27 & 0.28 & 0.31 & 0.48 & 0.65 & 0.72 & 0.64 \\
 & 0.05 & \textbf{0.76$^{+}$} & \textbf{0.76} & \textbf{0.76} & 0.58 & 0.66 & 0.66 & 0.44 & 0.64 & 0.29 & 0.41 & 0.10 & 0.22 & 0.25 & 0.27 & 0.35 & 0.55 & 0.51 & 0.49 \\
 & 0.10 & \textbf{0.71$^{+}$} & \textbf{0.71} & 0.70 & 0.49 & 0.53 & 0.65 & 0.38 & 0.50 & 0.26 & 0.32 & 0.08 & 0.18 & 0.22 & 0.23 & 0.28 & 0.43 & 0.42 & 0.43 \\
\hline
three big, many small & 0.00 & \textbf{0.83$^{+}$} & \textbf{0.83} & 0.82 & 0.76 & 0.81 & 0.72 & 0.33 & 0.40 & 0.21 & 0.27 & 0.07 & 0.18 & 0.15 & 0.18 & 0.47 & 0.64 & 0.74 & 0.66 \\
 & 0.05 & \textbf{0.78$^{+}$} & \textbf{0.78} & 0.74 & 0.48 & 0.67 & 0.68 & 0.28 & 0.34 & 0.16 & 0.22 & 0.06 & 0.13 & 0.14 & 0.16 & 0.35 & 0.51 & 0.42 & 0.44 \\
 & 0.10 & \textbf{0.67$^{+}$} & 0.66 & 0.62 & 0.36 & 0.55 & 0.61 & 0.23 & 0.29 & 0.14 & 0.17 & 0.05 & 0.10 & 0.12 & 0.14 & 0.29 & 0.43 & 0.33 & 0.40 \\
\hline
gradient & 0.00 & 0.42$^{-}$ & 0.40 & 0.39 & 0.36 & \textbf{0.61} & 0.56 & 0.21 & 0.21 & 0.11 & 0.14 & 0.03 & 0.14 & 0.07 & 0.16 & 0.50 & 0.48 & 0.41 & 0.38 \\
 & 0.05 & 0.32$^{-}$ & 0.31 & 0.27 & 0.20 & \textbf{0.48} & 0.40 & 0.14 & 0.15 & 0.07 & 0.10 & 0.03 & 0.07 & 0.07 & 0.11 & 0.34 & 0.42 & 0.18 & 0.16 \\
 & 0.10 & 0.23$^{-}$ & 0.24 & 0.22 & 0.13 & \textbf{0.39} & 0.35 & 0.09 & 0.11 & 0.06 & 0.07 & 0.02 & 0.05 & 0.06 & 0.09 & 0.27 & 0.33 & 0.12 & 0.13 \\
\hline
uniform $k{=}8$, $\alpha{=}1$ & 0.05 & \textbf{0.56$^{+}$} & 0.55 & 0.49 & 0.31 & 0.53 & 0.48 & 0.45 & 0.44 & 0.32 & 0.24 & 0.12 & 0.28 & 0.33 & 0.46 & 0.43 & 0.37 & 0.28 & 0.32 \\
uniform $k{=}8$, $\alpha{=}2$ & 0.05 & \textbf{0.56$^{+}$} & \textbf{0.56} & 0.49 & 0.33 & 0.54 & 0.48 & 0.45 & 0.45 & 0.32 & 0.22 & 0.12 & 0.28 & 0.33 & 0.46 & 0.43 & 0.39 & 0.30 & 0.32 \\
two worlds, $\alpha{=}1$ & 0.05 & \textbf{0.77$^{+}$} & \textbf{0.77} & 0.75 & 0.54 & 0.66 & 0.65 & 0.47 & 0.66 & 0.31 & 0.62 & 0.11 & 0.26 & 0.26 & 0.28 & 0.26 & 0.44 & 0.46 & 0.47 \\
two worlds, $\alpha{=}2$ & 0.05 & \textbf{0.79$^{+}$} & \textbf{0.79} & 0.75 & 0.51 & 0.66 & 0.64 & 0.49 & 0.66 & 0.32 & 0.62 & 0.10 & 0.29 & 0.26 & 0.30 & 0.23 & 0.41 & 0.50 & 0.50 \\
\hline
LFR, $o_n{=}0$ &  & 0.84$^{-}$ & 0.86 & 0.84 & 0.56 & 0.85 & 0.66 & 0.98 & 0.98 & \textbf{1.00} & 0.98 & 0.61 & \textbf{0.99} & 0.91 & \textbf{1.00} & 0.46 & 0.77 & 0.43 & 0.39 \\
LFR, $o_n{=}0.1$ &  & 0.69$^{-}$ & 0.72 & 0.71 & 0.27 & 0.53 & 0.60 & 0.78 & 0.78 & 0.88 & 0.81 & 0.49 & 0.87 & 0.79 & \textbf{0.94} & 0.41 & 0.72 & 0.29 & 0.23 \\
LFR, $o_n{=}0.3$ &  & 0.45$^{-}$ & 0.50 & 0.46 & 0.11 & 0.31 & 0.48 & 0.54 & 0.53 & 0.65 & 0.49 & 0.32 & 0.46 & 0.60 & \textbf{0.85} & 0.40 & 0.64 & 0.13 & 0.10 \\
LFR, $\langle k\rangle{=}10$ &  & 0.43$^{-}$ & 0.45 & 0.35 & 0.09 & 0.40 & 0.36 & 0.69 & 0.70 & 0.70 & 0.61 & 0.36 & 0.63 & 0.48 & \textbf{0.77} & 0.26 & 0.25 & 0.12 & 0.12 \\
LFR, $o_m{=}4$ &  & 0.61$^{-}$ & 0.63 & 0.62 & 0.40 & 0.46 & 0.52 & 0.62 & 0.62 & 0.70 & 0.62 & 0.40 & 0.67 & 0.68 & \textbf{0.75} & 0.43 & 0.66 & 0.36 & 0.28 \\
LFR, sizes 50--250 &  & 0.42$^{-}$ & 0.41 & 0.42 & 0.14 & 0.40 & 0.16 & 0.87 & 0.85 & 0.79 & 0.75 & 0.38 & 0.79 & 0.59 & \textbf{0.89} & 0.18 & 0.55 & 0.13 & 0.07 \\
LFR, $C_{\rm LFR}{=}0.3$ &  & 0.33$^{-}$ & 0.33 & 0.32 & 0.19 & 0.17 & 0.39 & 0.50 & 0.50 & 0.54 & 0.52 & 0.48 & 0.38 & \textbf{0.56} & 0.54 & 0.18 & 0.30 & 0.20 & 0.20 \\
LFR, $C_{\rm LFR}{=}0.5$ &  & 0.24$^{-}$ & 0.24 & 0.24 & 0.19 & 0.10 & 0.24 & 0.29 & 0.29 & 0.30 & 0.26 & \textbf{0.35} & 0.24 & \textbf{0.35} & 0.33 & 0.17 & 0.21 & 0.18 & 0.20 \\
\hline\hline
\end{tabular}}
\end{table*}


\emph{Label-free and oracle protocols.} Table~\ref{tab:labelfree} gives the label-free
$\Omega$ of every method and Table~\ref{tab:oracle} the oracle one. Under superlinear
preferential rewiring HyperCommon reaches $0.79$ on two worlds against $0.66$ for the
Constant Potts Model and link communities, and on uniform rings with $k=8$ it leads with
overlap ($0.50$ against $0.45$ for link communities at $o_n=0.10$). With labels, Leiden
with one objective or the other is ahead on most rings without overlap, on the dense
rings at $o_n=0.05$ and on rings of equal size and unequal degree; the Constant Potts
Model is ahead on two worlds at $o_n\le0.05$ by at most $0.05$. On the gradient at
$p=0$, the resolution at which Leiden returns the right number of communities gives
$\Omega=0.26$, and HyperCommon attains $0.97$ at the threshold at which it returns ten.
Under the native protocol HyperCommon is behind on the dense rings with little overlap,
where Leiden at $\gamma=1$ leads, and on the gradient; link communities lose most from
their own rule, the cut of maximum partition density returning $640$--$1320$
communities. The frozen and native tables are in the repository.

\emph{Per-community recovery and ONMI.} On the two families with a few large dense
rings and many small sparse ones, at the best threshold for the large rings part of each
small sparse ring has commonality below it, and HyperCommon splits such rings and leaves
part of their nodes unassigned (coverage $87\%$--$88\%$). By ONMI the Constant Potts
Model is then ahead on two worlds ($0.57$ against $0.42$) and on three big, many small
($0.52$ against $0.40$), where HyperCommon leads by $\Omega$. On the label-free covers,
with its resolution chosen by modularity, the Constant Potts Model recovers the small
rings less well and its ONMI lead on these two families is gone at $o_n=0.05$.

\subsection*{Robustness to noise}

\begin{table*}[t]
\caption{Robustness to noise on the ring families: mean $\Omega$ over the high-noise half ($p\in[0.4,0.8]$) and $p_{1/2}$, the rewired fraction at which the mean $\Omega$ over the ten test graphs first falls below $0.5$ (linear interpolation; $>0.8$: never), for HyperCommon (HC) and the best of the other methods by each measure, under the oracle and label-free protocols. Bold: HC ahead or within 0.005 (0.01 for $p_{1/2}$).}
\label{tab:robust}
\footnotesize\setlength{\tabcolsep}{3pt}
\begin{ruledtabular}
\begin{tabular}{lcccccccc}
 & \multicolumn{4}{c}{oracle} & \multicolumn{4}{c}{label-free} \\
Family & HC $\bar\Omega_{\rm hi}$ & best other & HC $p_{1/2}$ & best other & HC $\bar\Omega_{\rm hi}$ & best other & HC $p_{1/2}$ & best other \\
\hline
uniform, $k{=}8$, $o_n{=}0$ & 0.29 & 0.39 (SLPA) & 0.48 & 0.52 (Leiden) & 0.28 & 0.30 (SLPA) & \textbf{0.48} & 0.48 (L-CPM) \\
uniform, $k{=}8$, $o_n{=}0.05$ & 0.25 & 0.35 (SLPA) & \textbf{0.46} & 0.45 (Link) & 0.24 & 0.27 (SLPA) & \textbf{0.45} & 0.45 (Link) \\
uniform, $k{=}8$, $o_n{=}0.10$ & 0.22 & 0.31 (SLPA) & \textbf{0.42} & 0.42 (Link) & 0.21 & 0.24 (SLPA) & \textbf{0.41} & 0.38 (Link) \\
uniform, $k{=}16$, $o_n{=}0$ & 0.56 & 0.72 (Leiden) & 0.61 & 0.68 (Leiden) & 0.56 & 0.69 (Leiden) & 0.61 & 0.67 (Leiden) \\
uniform, $k{=}16$, $o_n{=}0.05$ & 0.47 & 0.59 (Leiden) & 0.58 & 0.64 (Leiden) & 0.46 & 0.56 (Leiden) & 0.58 & 0.63 (Leiden) \\
uniform, $k{=}16$, $o_n{=}0.10$ & 0.37 & 0.44 (Leiden) & 0.54 & 0.58 (Leiden) & 0.36 & 0.40 (Leiden) & 0.53 & 0.57 (Leiden) \\
\hline
equal $k$, unequal $s$ ($k{=}8$) & 0.19 & 0.26 (SLPA) & \textbf{0.41} & 0.40 (Link) & \textbf{0.19} & 0.19 (SLPA) & \textbf{0.41} & 0.40 (Link) \\
equal $k$, unequal $s$ ($k{=}16$) & 0.51 & 0.63 (Leiden) & 0.60 & 0.68 (Leiden) & 0.49 & 0.54 (Leiden) & 0.59 & 0.64 (Leiden) \\
equal $s$, unequal $k$ & 0.53 & 0.60 (L-CPM) & 0.66 & 0.69 (L-CPM) & 0.52 & 0.55 (L-CPM) & \textbf{0.66} & 0.66 (L-CPM) \\
two worlds & 0.61 & 0.62 (L-CPM) & \textbf{0.68} & 0.66 (Link) & \textbf{0.59} & 0.52 (Link) & \textbf{0.68} & 0.65 (Link) \\
three big, many small & \textbf{0.63} & 0.54 (Link) & \textbf{0.67} & 0.65 (Link) & \textbf{0.60} & 0.49 (Link) & \textbf{0.66} & 0.62 (Link) \\
gradient of ten rings & 0.23 & 0.26 (Link) & 0.47 & 0.48 (Link) & 0.14 & 0.18 (Link) & 0.17 & 0.41 (Link) \\
\hline
uniform $k{=}8$, pref.\ $\alpha{=}1$ & 0.25 & 0.35 (SLPA) & \textbf{0.46} & 0.46 (Leiden) & 0.24 & 0.27 (SLPA) & \textbf{0.46} & 0.45 (Link) \\
uniform $k{=}8$, pref.\ $\alpha{=}2$ & 0.25 & 0.34 (SLPA) & \textbf{0.46} & 0.46 (Leiden) & 0.24 & 0.27 (SLPA) & \textbf{0.46} & 0.45 (Link) \\
two worlds, pref.\ $\alpha{=}1$ & 0.60 & 0.70 (L-CPM) & 0.67 & 0.70 (L-CPM) & \textbf{0.59} & 0.50 (L-CPM) & 0.67 & 0.69 (L-CPM) \\
two worlds, pref.\ $\alpha{=}2$ & 0.61 & 0.64 (L-CPM) & \textbf{0.67} & 0.66 (L-CPM) & \textbf{0.60} & 0.48 (L-CPM) & \textbf{0.67} & 0.62 (L-CPM) \\
\end{tabular}
\end{ruledtabular}
\end{table*}


Table~\ref{tab:robust} reports, for every ring family, the rewired fraction $p_{1/2}$ at
which the mean $\Omega$ first falls below one half and the mean $\Omega$ over the
high-noise half, $p\in[0.4,0.8]$. Without labels, on the four families whose large rings
have degree 32 and small rings degree 8, HyperCommon has the highest high-noise mean
($0.60$ against $0.49$ for link communities on three big, many small); on the gradient
it is the least robust of the leading methods. Only DEMON and ANGEL degrade under
degree-biased noise; for HyperCommon the bound of Jaccard by the degree ratio offers a
mechanism: a hub grown by preferential rewiring has low commonality with the lattice
neighbours it attaches to, so the edges it gains rarely pass.

\subsection*{Ablations and the commonality function}

Without the third-pair test, a wedge whose two edges pass reaches into another community
through a shared node or through an edge between communities, which is how the two-test
variant loses $0.21$ even on LFR graphs without overlap. Triangles alone are better than
HyperCommon in two settings under the oracle protocol and in none under the label-free
one; with one parameter frozen per family they have the higher mean $\Omega$ in 15 of
the 36 settings. Joining triples only through shared edges is slightly better on the LFR
graphs, by up to $0.03$ with labels and $0.05$ without.

Tables~\ref{tab:familyoracle} and~\ref{tab:familyfree} run HyperCommon with each
commonality function on eighteen settings (six uniform-ring, six unequal-ring at
$o_n=0.05$ and six LFR), the threshold taken as a quantile of each function's own triple
scores. With the best threshold Jaccard is the best, or within $0.01$, on seventeen and
never more than $0.02$ behind. The two unnormalised scores, the raw overlap and resource
allocation, are less accurate than Jaccard on all eighteen, by up to $0.21$ and $0.20$:
where communities differ in degree, a threshold on an unnormalised score is a threshold
on degree as much as on commonality. Choosing the function together with the threshold
raises $\Omega$ by $0.01$--$0.06$. Without labels, on the gradient the selected cover of
$I^2/U$ reaches $\Omega=0.50$ against $0.36$ for Jaccard, but the selector, asked to
choose the function as well, does not find it ($0.31$).

\subsection*{LFR graphs, geometry, controls and null graphs}

With its best threshold HyperCommon reaches $\Omega=0.885$ at $\mu=0.1$ and $o_n=0.1$,
against $0.884$ for Leiden with modularity and for Infomap and $0.886$ for Leiden with
the Constant Potts Model, while SLPA reaches $0.99$ [Fig.~\ref{fig:benchmarksoracle}(f)].
At $\mu=0.3$ only $5\%$ of the wedges inside communities pass at $t=0.1$ in the slice with
communities of 50 to 250 nodes. The generator's own rewiring, used to raise the
clustering, adds triangles across communities as well: the share of triangles spanning
two communities grows from $2\%$ to $19\%$ and $32\%$.

In the geometry families the mean clustering of the noiseless graphs falls from
$0.62$--$0.68$ at $q=0$ to $0.48$--$0.51$ at $q=0.1$ and $0.04$--$0.16$ at $q=1$
(Table~\ref{tab:geometry}). On lattice communities HyperCommon leads with labels by
$0.94$ against $0.87$ for clique percolation on three big, many small and is level with
the Constant Potts Model on two worlds ($0.91$); with labels at $q=1$ on the uniform rings
it reaches $0.60$ against $0.85$. Open wedges slow the decline: there triangles alone fall
from $0.65$ to $0.29$ as $q$ goes from 0 to 1, and HyperCommon from $0.73$ to $0.60$.

Among the controls (Table~\ref{tab:S-controls}), with hubs inside the communities the
Leiden objectives gain and lead, Leiden on the dense hub rings ($0.83$ against $0.69$
without labels) and the Constant Potts Model on two worlds with hubs ($0.83$ against
$0.75$). With split overlap HyperCommon is level with the best other method or ahead on
two worlds and three big, many small under both protocols. With random geometric
communities it is behind Leiden with labels ($0.50$ against $0.57$) and only slightly
ahead of it without labels on the uniform rings ($0.48$ against $0.46$), and behind the
Constant Potts Model on two worlds under both protocols ($0.73$ against $0.84$ and
$0.79$). On the null graphs (Table~\ref{tab:S-nulls}) label-free HyperCommon covers
$98\%$ and $96\%$ of the nodes with its $101$ and $175$ communities, against $32$--$36$
and $12$--$28$ for Leiden, Infomap, OSLOM and the stochastic block model.

\subsection*{Real networks}

On the SNAP graphs, Leiden's label-free F1 equals its F1 at $\gamma=1$
(Table~\ref{tab:real}). With labels Infomap leads by F1 on Amazon ($0.935$), OSLOM on
DBLP ($0.549$, against $0.545$ for Infomap) and Leiden on YouTube ($0.363$); Infomap at
its default Markov time is about $0.06$ ahead of label-free HyperCommon on Amazon and
DBLP. On Amazon, whose communities are sparse in triangles, the two-test variant reaches
F1 $0.887$ but a lower $\Omega$ ($0.424$ against $0.459$). Detecting on the subgraphs
induced by the labelled nodes, as is common, makes F1 uninformative on Amazon and
YouTube: there the connected components score within $0.005$ of the best method.

On the fourteen smaller networks (Tables~\ref{tab:realmore}, \ref{tab:realmoreomega},
\ref{tab:S-realnet-f1-best} and~\ref{tab:S-realnet-f1-free}), HC$_f$ chooses, with
labels, among 700 (function, threshold) pairs per network, an upper bound rather than one
parameter grid; it is the most accurate method on four networks by F1 and four by
$\Omega$. OSLOM has the best mean rank without labels ($6.2$) and Leiden with labels
($6.0$); OSLOM, Leiden, SLPA, label propagation, SCAN and Infomap each lead on some
networks. The label-free threshold retains $95\%$ of the best F1 or more on seven of the
fourteen; on political blogs, a two-group network in which HyperCommon's best cover is the
giant component, it selects fifteen small communities instead.
